


\section{Représentation de documents}

\begin{frame}
\frametitle{Représentation de documents textuels}

La représentation des données est basée sur des \alert{graphes}, permettant de représenter des
données régulières ou irrégulières.

\vfill

\alert{Notions essentielles}

\begin{description}
  \item[Les données sont auto-décrites.] Le contenu vient avec sa propre
  description.
  
 \item[Structures riches.] Le contenu se décrit avec des listes, des
 enregistrements imbriqués, des ensembles.

  \item[Typage flexible.] Les données peuvent être typées (``c'est un entier''),
  et/ou structurées (``cette partie du graphe \alert{doit} avoir telle forme''), et/ou ni l'un ni
  l'autre. Dans ce dernier cas c'est à l'application d'effectuer le contrôle.
 
   \item[Sérialisation.] Structure \alert{et} contenu doivent pouvoir être transformés
      ensemble en une chaîne de caractères autonome.

\end{description}

\end{frame}

\begin{frame}[fragile]
\frametitle{Commençons avec JSON}

Point de départ\,: \alert{listes associatives}, ensemble de paires (clé, valeur).

\begin{minted}{javascript}
{"nom": "Philippe", 
 "tel": 2157786,  
 "email": "philippe@cnam.fr"}
\end{minted}

Extension naturelle\,: les valeurs sont elles-mêmes d'autres structures.

\begin{minted}{javascript}
{"nom": {"prenom": "Philippe",  "famille": "Rigaux"},
 "tel": 2157786,
 "email": "philippe@cnam.fr"}
\end{minted}

Autre extension: imbrication de listes.

\begin{minted}{javascript}
    {"nom": "Philippe", "tels": [2157786, 2498762] }
\end{minted}

\begin{remblock}{Remarque}
Les exemples ci-dessus montrent un \alert{codage de l'information} sous forme
d'une chaîne 
de caractères stockable sur disque et échangeable par réseau.
\end{remblock}
\end{frame}

\begin{frame}
\frametitle{Vue conceptuelle: modélisation  arborescente}

Le codage représente un \alert{arbre}. Les étiquettes sont
sur les arêtes, les valeurs sur les feuilles.


\includegraphics[width=11cm]{label_on_edges1}

\begin{remblock}{Important}
Le codage distingue bien une \alert{structure} (l'arbre) et un 
\alert{contenu} (le texte dans les feuilles). Toutes les  \alert{opérations} (de recherche, de navigation, ...)
s'appuient sur la structure.
\end{remblock}
\end{frame}

\begin{frame}
\frametitle{Autre possibilité\,: les étiquettes sont des  n{\oe}uds}

On peut représenter à la fois la structure et les valeurs 
comme des n{\oe}uds.

\medskip

\includegraphics[width=9cm]{simple_xml_tree}

\medskip

\begin{remblock}{Remarque}
C'est le choix de représentation du langage XML: toute l'information (structure et contenu)
est stockée sous forme de n{\oe}ud.
\end{remblock}

\end{frame}

\begin{frame}{JSON, Javascript Object Notation}

\begin{redbullet}
\item JavaScript Object Notation;
\item Initialement créé pour la sérialisation et l'échange d'objets JavaScript;
\item Langage pour l'échange de données semi-structurées (et éventuellement structurées);
\item Format texte indépendant du language de programmation utilisé pour le manipuler.
\end{redbullet}

\vfill 

Utilisation première : échange de données dans un environnement Web (par exemple applications Ajax)\\

\vfill

Extension : sérialisation
 et stockage de données 

\end{frame}

\definecolor{lightgray}{rgb}{.9,.9,.9}
\definecolor{darkgray}{rgb}{.4,.4,.4}
\definecolor{purple}{rgb}{0.65, 0.12, 0.82}


\begin{frame}[containsverbatim,fragile,allowframebreaks]{Les bases de JSON}

Structure de base: paire  \textit{clef-valeur} (\textit{key-value})


\begin{minted}{javascript}
"title": "The Social network"
\end{minted}

\bigskip

\alert{Qu'est-ce qu'une valeur?} On distingue les
 valeurs \alert{atomiques} et les valeurs \alert{complexes} (construites)

\medskip
Valeurs atomiques\,: cha\^ines de caractères (entourées 
par les classiques guillemets anglais (droits)), nombres (entiers, flottants) 
et valeurs booléennes (\texttt{true} ou \texttt{false}).

\begin{minted}{javascript}
"year": 2010
\end{minted}


\begin{minted}{javascript}
"oscar": false
\end{minted}

\framebreak

Valeurs complexes\,: les \alert{objets}.
Un \textit{objet} est un ensemble de paires clef-valeur.

\vfill

Au sein d'un ensemble de paires, une clef apparait au plus une fois 
(NB : les types de valeurs peuvent être distincts).

\begin{minted}{javascript}
{"last_name": "Fincher", "first_name": "David"}
\end{minted}

\vfill

Un objet peut être utilisé comme valeur (dite \textit{complexe}) dans une paire clef-valeur.

\begin{minted}{javascript}
"director": { 
  "last_name": "Fincher",
  "first_name": "David", 
  "birth_date": 1962
}
\end{minted}

\framebreak

Valeurs complexes\,: les \alert{tableaux}.

\bigskip

Un \textit{tableau} (\textit{array}) est une liste 
 de valeurs (dont le type n'est pas forcément le même).

\begin{minted}{javascript}
"actors": ["Eisenberg", "Mara", "Garfield", "Timberlake"]
\end{minted}

\bigskip

Imbrication sans limite\,: tableaux de tableaux, tableaux d'objets contenant eux-mêmes des tableaux, etc.
\framebreak

Un \textit{document} est un objet. Il peut être défini par des objets et tableaux imbriqu\' es autant de fois que nécessaire.

\begin{minted}[baselinestretch=.9,
               fontsize=\small,
               framesep=2mm]{javascript}
{
  "title": "The Social network", 
  "summary": "On a fall night in 2003, Harvard undergrad and \n
         programming genius Mark Zuckerberg sits down at his \n 
         computer and heatedly begins working on a new idea. (...)",
  "year": 2010, 
  "director": {"last_name": "Fincher",
                    "first_name": "David"},
  "actors": [ 
        {"first_name": "Jesse", "last_name": "Eisenberg"}, 
        {"first_name": "Rooney", "last_name": "Mara"}
    ]
}
\end{minted}

\end{frame}



\begin{frame}[fragile]
\frametitle{Flexibilité de la représentation\,: un rapport écrit}

\begin{minted}[baselinestretch=.9,
               fontsize=\small,
               framesep=2mm]{javascript}
{
    "author": "Philippe Rigaux"
    "date": "04/10/2065",
    "title": "BD documentaires",
    "content": {
       "para": "...",
       "para": "...",
       "section": {
          "title": "...",
          "para": "...",
          "figure" {"caption": "...", "content:"},
          "para": "...",
          "section": {...}
       }
    }
}
\end{minted}

\begin{redbullet}
  \item Imbrication très libre d'éléments représentant des composants d'un rapport.
  \item Pas ou peu de \alert{référence} (mais on pourrait l'envisager pour l'auteur).
\end{redbullet}
\end{frame}




\begin{frame}[fragile]
\frametitle{Ce qu'il faut retenir}

Les documents semi-structurés se caractérisent par les aspects suivants.

\begin{redbullet}
  \item Une structure \alert{flexible} (présence/absence, variations), \alert{complexe} (imbrication),
   et \alert{auto-décrite}.
  \item Une double représentation: un vue conceptuelle (arbre) et son codage (forme sérialisée).
    \item \alert{Valable} pour les documents "multimédia" au sens large: rapports, images, vidéos.
 \end{redbullet}
 
\vfill
La notion d'autonomie (pas d'association entre documents) est essentielle.
\begin{redbullet}
       \item Autonomie? On peut mettre en {\oe}uvre un \alert{passage à l'échelle par
       distribution}. Typique des systèmes NoSQL.
       \item Pas d'autonomie: un jour ou l'autre il faudra faire des jointures; 
          les systèmes NoSQL ne sont pas bon pour ça.
 \end{redbullet}

\end{frame}

